% 【本文件写什么】impl 实现层：kernel
% - 定位：内核内建伪 FS / rootfs 逻辑，devfs、procfs、rootfs。
% - crate 路径：os/components/wateros-syscall/syscall-impl/impl-kernel/src/
% - 如何实现对应 api-v0 trait：关键类型、状态机、数据结构
% - 算法或硬件交互流程，可用伪代码/流程图
% - 本 impl 启用的 Cargo [features] 与 arch feature，impl-riscv64 等
% - 与其它 impl 变体的差异、切换 feature 时的影响
% - 性能/内存假设与已知限制
% 【事实来源】
% - 上述 crate 源码与 Cargo.toml
% - os/feature-tree.txt 中指向本 impl 的 feature 链

\paragraph{分发架构}

\code{KernelSyscallDispatcher}实现\code{SyscallDispatcher}，关联\code{ActiveSyscallNumberTable}，主线为 \code{LinuxGeneric64}。热路径\code{dispatch\_syscall\_by\_nr}，见 \code{syscall\_nr\_dispatch.rs}，对号表收录的裸号做单次\code{match}，直接调用对应\code{dispatch\_*}；未命中时先走\code{dispatch\_syscall\_aliases}处理旁路号，再返回\code{-ENOSYS}。

旁路号包括\code{statx}，291、\code{fstatat}，79、\code{faccessat2}，439、\code{close\_range}，436、\code{setgroups}，159、\code{sched\_setattr/getattr}、\code{adjtimex}/\code{clock\_adjtime}、\code{acct}，以及xattr裸号5--16，\code{setxattr} 至 \code{fremovexattr}。这些号在\code{LinuxGeneric64}表中有定义但分发路径单独处理，避免与主\code{match}重复。

各\code{dispatch\_*}内联转发到\code{sys::sys\_*}，返回值取\code{UserRet}编码后的\code{.0}。未在\code{impl-kernel}覆盖的槽位覆盖\code{dispatch\_unsupported}为\code{panic}，bring-up策略，与trait默认的\code{-ENOSYS}不同。

\paragraph{sys 模块划分}

\code{sys/mod.rs}按域组织实现文件，代表性覆盖包括：
\begin{itemize}
 \item 文件与目录：\code{read}、\code{write}、\code{openat}、\code{close}、\code{lseek}、\code{fcntl}、\code{getdents64}、\code{mount}、\code{umount2}、openat/linkat等at族syscall、xattr、\code{stat}/\code{fstat}/\code{statx}、截断与\code{fallocate}；
 \item 内存：\code{brk}、\code{mmap}族、\code{mprotect}、\code{mremap}、\code{madvise}、\code{mlock*}、\code{get\_mempolicy}、SysV\code{shm}子集；
 \item 进程/线程：\code{clone}/\code{clone3}、\code{execve}、\code{exit}/\code{exit\_group}、\code{waitpid}/\code{waitid}、\code{getpid}族、\code{sched\_*}、\code{unshare}，含 \code{CLONE\_NEWNS}，robust futex列表；
 \item 凭证：\code{cred.rs}经\code{wateros-cred}门面处理\code{getuid}族与\code{set*id}；\code{cap.rs}提供\code{capget}/\code{capset}最小桩；
 \item 信号/同步：\code{futex}、\code{rt\_sig*}、\code{kill}/\code{tkill}/\code{tgkill}、\code{nanosleep}/\code{clock\_nanosleep}；
 \item 多路复用：\code{poll}/\code{ppoll}/\code{select}/\code{pselect6}、\code{epoll\_*}，\code{poll\_engine} 与 \code{epoll\_fd} 协作；
 \item 网络：INET socket全族与\code{AF\_UNIX}，含 \code{socket\_fd}、\code{socket\_block}、\code{unix\_sock}；
 \item 杂项：\code{syslog}转发\code{klog}、\code{sysinfo}、\code{uname}、\code{prctl}子集、\code{getrandom}、LTP fast-exit协作，含 \code{ltp\_cgroup\_helper}、\code{bringup\_stats}。
\end{itemize}

\paragraph{trap协作}

除syscall分发外，本crate导出：
\begin{itemize}
 \item \code{timer\_tick(interrupted\_user)}：时钟tick，协作信号与调度；
 \item \code{deliver\_pending\_signal}、\code{restore\_signal\_frame}：信号帧投递与恢复；
 \item \code{raise\_current\_signal}：向当前任务发信号；
 \item \code{is\_restartable\_syscall\_nr}：EINTR后可自动重启的号查询，O，1跳表；
 \item \code{drop\_reaped\_task\_runtime\_resources}：\code{waitpid} reap后清理syscall侧资源，fd、epoll、socket等；
 \item \code{record\_user\_page\_fault\_handled}、\code{log\_thread\_bringup\_stats\_summary}：bring-up统计。
\end{itemize}

\paragraph{用户内存与VFS辅助}

\code{user\_copy}、\code{fallible\_buf}集中处理\code{copy\_from\_user}/\code{copy\_to\_user}与可失败缓冲区；\code{vfs\_util}、\code{mm\_util}、\code{linux\_stat}、\code{stat\_times}提供路径解析、统计结构填充等共用逻辑。VFS、MM、task、ipc、cred各组件只暴露内核侧稳定接口，syscall层负责Linux ABI到这些接口的转换与\code{EFAULT}/\code{EINVAL}等错误映射。

\paragraph{已知限制}

\code{utimensat}、部分\code{ioctl}/控制终端、真实权限校验仍是最小语义。\code{SyscallKind}齐全不等于\code{sys\_*}已实现；以本impl源码为准。RISC-V64与LoongArch64共用generic 64位号表，架构差异留在\code{wateros-platform} trap层。
